\documentclass[11pt]{article}
\usepackage[margin=1in]{geometry}
\usepackage{amsmath,amssymb,mathtools}
\usepackage{booktabs,longtable}
\usepackage{enumitem}
\usepackage[hidelinks]{hyperref}

\newcommand{\dd}{\mathrm{d}}
\newcommand{\R}{\mathbb{R}}
\newcommand{\TT}{\mathrm{TT}}
\newcommand{\calI}{\mathcal{I}}
\newcommand{\calB}{\mathcal{B}}
\newcommand{\calM}{\mathcal{M}}
\newcommand{\dotp}{\partial_i}

\title{Fixed-Branch Analytical Derivative for the Zhao--Cui Tensor-Train Sequential Filter}
\author{BayesFilter P11 Derivation Note}
\date{2026-05-30}

\begin{document}
\maketitle

\begin{abstract}
This note derives the analytical derivative of the scalar produced by a
fixed-branch Zhao--Cui squared tensor-train sequential filter.  The derivative
is for the declared approximate scalar
\[
  \widehat\ell_T^{\TT}(\alpha)
  =
  \sum_{t=1}^T \{\log \widehat z_t(\alpha)-c_t(\alpha)\},
\]
where \(\alpha\) is an external model parameter, \(\widehat z_t\) is the
squared-TT/SIRT normalizer, and \(c_t\) is the stabilizing constant used in the
companion code's \texttt{log(sirt.z)-const} accumulator.  The derivation is
complete only for a frozen algorithmic branch: fixed variable ordering, fixed
bases, fixed domains, fixed interpolation or cross sets, fixed ranks, fixed
rounding decisions, fixed random/debug samples, fixed defensive term policy,
fixed preconditioner choice, and fixed selected minimizer for \(c_t\).  It does
not claim a complete derivative for the adaptive stochastic companion code.
\end{abstract}

\section{Anchors And Scope}

The source construction is Zhao--Cui's sequential TT method.  Their Algorithm 1
forms an unnormalized density
\[
  q_t(x_t,\vartheta,x_{t-1})
  =
  \widehat\pi_{t-1}(x_{t-1},\vartheta\mid y_{1:t-1})
  f(x_t\mid x_{t-1},\vartheta)
  g(y_t\mid x_t,\vartheta),
  \tag{ZC-12}
\]
then reapproximates it by a TT and integrates the TT to obtain a normalizing
constant.  Their squared-TT construction approximates the square root of an
unnormalized density by
\[
  \sqrt{\pi(x)}
  \approx
  \phi(x)=H_1(x_1)\cdots H_m(x_m),
\]
and defines the nonnegative defensive approximation
\[
  \widehat\pi(x)=\phi(x)^2+\tau\lambda(x),\qquad
  \widehat z=\int \widehat\pi(x)\,\dd x.
  \tag{ZC-13}
\]
Algorithm 2 applies this construction to the sequential density \(q_t\), using
the square-root TT \(\phi_t\) and defensive density
\[
  \widehat\pi_t(x_t,\vartheta,x_{t-1})
  =
  \phi_t(x_t,\vartheta,x_{t-1})^2
  +
  \tau_t\lambda(x_t)\lambda(\vartheta)\lambda(x_{t-1}).
  \tag{ZC-16}
\]
Zhao--Cui Section 4.1 identifies the exact joint normalizer as evidence and
explains that the algorithm approximates \(q_t\) because the exact marginal
posterior entering the recursion is not directly evaluable.

The inspected companion code implements the same scalar path through:
\[
\begin{array}{ll}
\texttt{models/full\_sol.m} &
  \texttt{sol.logmarginal\_likelihood += log(sirt.z)-const},\\
\texttt{deep-tensor.dev/src/SIRT.m} &
  \texttt{potential\_to\_density},\\
\texttt{deep-tensor.dev/src/@TTSIRT/TTSIRT.m} &
  \texttt{approx = TTFun(func,...)} ,\\
\texttt{deep-tensor.dev/src/@TTSIRT/marginalise.m} &
  \texttt{obj.z = obj.fun\_z + obj.tau}.
\end{array}
\]
The P10 Octave audit snapshot under
\texttt{third\_party/audit/zhao\_cui\_tensor\_ssm\_p10/source} shows that a
reduced Kalman path reaches the TT/SIRT approximation, completes one ALS pass,
and reports a finite log-marginal-likelihood accumulator.  That smoke is
implementation evidence, not a proof of posterior accuracy.

\paragraph{Parameter notation.}
Zhao--Cui use \(\theta\) for the learned state-space parameter inside the
sequential posterior.  This note writes that variable as \(\vartheta\).  The
external parameter with respect to which we differentiate the approximate
scalar is \(\alpha_i\), and \(\dot h\) means \(\partial_i h\).  If a future
BayesFilter implementation samples \(\vartheta\) directly by HMC, it must first
translate this external-parameter derivative into the chosen target
parameterization.

\section{Declared Scalar}

Let
\[
  u_t=(x_t,\vartheta,x_{t-1})\in\R^D.
\]
At time \(t\), the unnormalized density supplied to the squared-TT step is
\[
  q_t(u_t;\alpha)
  =
  \widehat p_{t-1}^{\TT}(x_{t-1},\vartheta;\alpha)
  f_\alpha(x_t\mid x_{t-1},\vartheta)
  g_\alpha(y_t\mid x_t,\vartheta).
  \label{eq:p11-qt}
\]
The previous approximation \(\widehat p_{t-1}^{\TT}\) is the marginal of the
previous squared-TT density after normalization.  At \(t=1\), it is replaced by
the prior and initial-state density used by the model.

The code constructs an affine approximation coordinate
\[
  u_t = A_t r_t+b_t,
  \label{eq:p11-affine}
\]
with \(A_t=\texttt{L\_temp}\) and \(b_t=\texttt{mu\_temp}\).  In the companion
code, \(A_t\) and \(b_t\) are estimated from weighted samples and are therefore
branch-changing unless frozen.  In this note they are fixed smooth branch
objects.  The unshifted potential in \(r\)-coordinates is
\[
  V_{t,0}(r;\alpha)
  =
  -\log q_t(A_t r+b_t;\alpha)-\log|\det A_t|.
  \label{eq:p11-v0}
\]
The code then chooses a stabilizing constant
\[
  c_t(\alpha)=V_{t,0}(r_{t,s^\star};\alpha)
  \label{eq:p11-const}
\]
when the selected minimizer \(s^\star\) of the fixed random/debug grid is
frozen and unique.  If the implementation freezes the numerical value of
\(c_t\) instead, then \(\dot c_t=0\), but the differentiated scalar is a
different declared scalar.  The shifted potential is
\[
  V_t(r;\alpha)=V_{t,0}(r;\alpha)-c_t(\alpha).
  \label{eq:p11-v}
\]

The SIRT implementation approximates a square-root density in the reference
coordinates used by the basis.  Let \(\rho_t(r;\alpha)\) denote the
unnormalized reference-coordinate density after all fixed basis-domain
Jacobian and reference-measure factors used by \texttt{SIRT.potential\_to\_density}
have been included.  Then
\[
  \psi_t(r;\alpha)=\rho_t(r;\alpha)^{1/2}
  \label{eq:p11-psi}
\]
is the function approximated by \texttt{TTFun}.  In the simplest affine and
fixed-basis case,
\[
  \dot\psi_t(r;\alpha)
  =
  -\frac12\,\psi_t(r;\alpha)\,\dot V_t(r;\alpha),
  \label{eq:p11-psi-derivative}
\]
with additional Jacobian or reference-density derivative terms only if the
basis-domain map itself depends on \(\alpha\).

The fixed-branch squared-TT approximation is
\[
  \phi_t(r;\alpha)
  =
  G_{t,1}(r_1;\alpha)\,
  G_{t,2}(r_2;\alpha)\cdots
  G_{t,D}(r_D;\alpha),
  \label{eq:p11-tt-function}
\]
where \(G_{t,k}(r_k;\alpha)\in\R^{r_{k-1}\times r_k}\),
\(r_0=r_D=1\), all ranks are fixed, and every basis/interpolation decision is
fixed.  The defensive squared-TT normalizer is
\[
  \widehat z_t(\alpha)
  =
  \int \phi_t(r;\alpha)^2\,\dd \nu(r)
  +
  \tau_t(\alpha),
  \label{eq:p11-z}
\]
where \(\nu\) is the fixed tensor-product basis/reference measure and
\(\tau_t\) is the defensive mass term.  Thus the declared scalar is
\[
  \widehat\ell_T^{\TT}(\alpha)
  =
  \sum_{t=1}^T
  \{\log \widehat z_t(\alpha)-c_t(\alpha)\}.
  \label{eq:p11-scalar}
\]

\paragraph{Regularity assumptions.}
The formulas below require the usual finite-dimensional sensitivity
conditions on the frozen branch.  The transition and observation densities,
the previous TT marginal, and the fixed coordinate maps are differentiable in
\(\alpha\) on the evaluated domain; the normalizers \(R_{t,0}+\tau_t\) are
finite and positive; the support does not change pathologically with
\(\alpha\); the fixed local linear or least-squares systems are nonsingular
or have a declared differentiable pseudoinverse branch; and differentiation
may be interchanged with the finite contractions or with the stated integrals.
If QR, SVD, Cholesky, or inverse-CDF solves are differentiated, their active
rank, sign, bracket, and simple-spectrum branches are part of the fixed branch.

\section{Derivative Of The Input Density}

From \eqref{eq:p11-qt},
\[
  \dot{\log q_t}
  =
  \dot{\log \widehat p_{t-1}^{\TT}}
  +
  \dot{\log f_\alpha(x_t\mid x_{t-1},\vartheta)}
  +
  \dot{\log g_\alpha(y_t\mid x_t,\vartheta)}.
  \label{eq:p11-logq-derivative}
\]
The model-score terms are ordinary transition and likelihood derivatives.  The
recursive term is more delicate.  If
\[
  \widehat p_{t-1}^{\TT}(v;\alpha)
  =
  \frac{a_{t-1}(v;\alpha)}{\widehat z_{t-1}(\alpha)}
  \]
is the marginal density over \(v=(x_{t-1},\vartheta)\), then
\[
  \dot{\log \widehat p_{t-1}^{\TT}}(v;\alpha)
  =
  \frac{\dot a_{t-1}(v;\alpha)}{a_{t-1}(v;\alpha)}
  -
  \frac{\dot{\widehat z}_{t-1}(\alpha)}
        {\widehat z_{t-1}(\alpha)}.
  \label{eq:p11-prev-score}
\]
The marginal numerator \(a_{t-1}\) is computed by contracting the previous
TT cores over the eliminated coordinates.  Its derivative is obtained by the
same contraction with one differentiated core at a time, plus any derivative
of \(\tau_{t-1}\lambda\) if the defensive term or reference density depends on
\(\alpha\).

If \(A_t\), \(b_t\), or the selected constant point \(r_{t,s^\star}\) depend
smoothly on \(\alpha\), then
\[
\begin{aligned}
  \dot V_{t,0}(r)
  &=
  -\dot{\log q_t}(u_t)
  -
  \nabla_{u}\log q_t(u_t)^\top
      (\dot A_t r+\dot b_t)
  -
  \operatorname{tr}(A_t^{-1}\dot A_t),\\
  \dot c_t
  &=
  \dot V_{t,0}(r_{t,s^\star})
  +
  \nabla_r V_{t,0}(r_{t,s^\star})^\top \dot r_{t,s^\star}.
  \label{eq:p11-v-derivative}
\end{aligned}
\]
For the companion code's random grid, \(r_{t,s^\star}\) is fixed after branch
freezing, so \(\dot r_{t,s^\star}=0\).  If the minimizing index changes, the
ordinary derivative is not the derivative of a single smooth branch.

\section{Fixed-Branch TT Core Sensitivities}

Equation \eqref{eq:p11-tt-function} gives the product-rule derivative
\[
  \dot\phi_t(r)
  =
  \sum_{k=1}^D
  G_{t,1}(r_1)\cdots
  \dot G_{t,k}(r_k)\cdots
  G_{t,D}(r_D).
  \label{eq:p11-tt-product-rule}
\]
The missing implementation object is therefore \(\dot G_{t,k}\) for each
fixed core.  A fixed TT-cross or ALS branch can be written as a finite system
of interpolation or least-squares equations.  Let \(g_{t,k}\) be the vector of
free coefficients in core \(k\), with all ranks, basis functions, sample
points, sweep order, and neighboring interfaces fixed.

For an exact fixed interpolation equation,
\[
  A_{t,k}(\alpha)g_{t,k}(\alpha)=b_{t,k}(\alpha),
  \label{eq:p11-core-linear}
\]
the sensitivity is
\[
  A_{t,k}\dot g_{t,k}
  =
  \dot b_{t,k}-\dot A_{t,k}g_{t,k}.
  \label{eq:p11-core-linear-derivative}
\]
Here \(b_{t,k}\) consists of target square-root evaluations
\(\psi_t(r^{(j)};\alpha)\) after subtracting or conditioning on the fixed
left/right interfaces, and \(\dot b_{t,k}\) uses
\eqref{eq:p11-psi-derivative}.  The matrix \(A_{t,k}\) contains basis
evaluations and interface products; \(\dot A_{t,k}\) is zero only when the
interfaces are held numerically fixed.  If interfaces from earlier core
updates are differentiated, \(\dot A_{t,k}\) contains their sensitivities.

For a weighted least-squares core update,
\[
  g_{t,k}
  =
  \arg\min_g \|W_{t,k}^{1/2}(A_{t,k}g-b_{t,k})\|_2^2,
  \label{eq:p11-core-ls}
\]
with fixed weights and full-rank normal matrix
\(N_{t,k}=A_{t,k}^\top W_{t,k}A_{t,k}\), the normal-equation derivative is
\[
\begin{aligned}
  N_{t,k}\dot g_{t,k}
  &=
  \dot A_{t,k}^\top W_{t,k}(b_{t,k}-A_{t,k}g_{t,k})
  +A_{t,k}^\top W_{t,k}(\dot b_{t,k}-\dot A_{t,k}g_{t,k}).
  \label{eq:p11-core-ls-derivative}
\end{aligned}
\]
If the residual vanishes, this reduces to the interpolation formula.  If rank
truncation, pivot selection, enrichment samples, or ALS stopping change, then
\eqref{eq:p11-core-linear-derivative}--\eqref{eq:p11-core-ls-derivative}
describe only the derivative of the frozen local branch.

\section{Normalizer Contraction And Its Derivative}

Let \(\{b_{k,a}\}_{a=1}^{n_k}\) be the fixed one-dimensional basis on
coordinate \(k\), and let
\[
  M_k[a,b]=\int b_{k,a}(r_k)b_{k,b}(r_k)\,\dd\nu_k(r_k)
  \label{eq:p11-mass}
\]
be the one-dimensional mass matrix.  Write the coefficient core slices as
\[
  C_{t,k,a}(\alpha)\in\R^{r_{k-1}\times r_k}.
\]
Then
\[
  G_{t,k}(r_k;\alpha)=\sum_{a=1}^{n_k} C_{t,k,a}(\alpha)b_{k,a}(r_k).
\]
Define the right contractions
\[
  R_{t,D}=1,\qquad
  R_{t,k-1}
  =
  \sum_{a,b=1}^{n_k}
  M_k[a,b]\,
  C_{t,k,a} R_{t,k} C_{t,k,b}^\top,
  \quad k=D,\ldots,1.
  \label{eq:p11-right-contraction}
\]
The square part of the normalizer is the scalar
\[
  z_{t,\phi}=R_{t,0}.
  \label{eq:p11-zphi}
\]
This is the direct mass-matrix form of the normalizer.  The companion code's
\texttt{@TTSIRT/marginalise.m} computes the same quantity in exact arithmetic
through repeated mass-weighted QR contractions and then sets
\[
  \texttt{sirt.z}=\texttt{sirt.fun\_z}+\texttt{sirt.tau}.
\]
For a strict same-code-scalar derivative, one may differentiate the QR branch.
For a cleaner implementation, compute both value and derivative by the direct
contraction \eqref{eq:p11-right-contraction}; the scalar must then be declared
as that contraction rather than as a black-box QR floating-point value.

Because the basis is fixed in this note, \(\dot M_k=0\).  The contraction
derivative is
\[
\begin{aligned}
  \dot R_{t,k-1}
  =
  \sum_{a,b=1}^{n_k}M_k[a,b]\Big[
  &\dot C_{t,k,a}R_{t,k}C_{t,k,b}^\top
  +C_{t,k,a}\dot R_{t,k}C_{t,k,b}^\top \\
  &+C_{t,k,a}R_{t,k}\dot C_{t,k,b}^\top
  \Big],
  \qquad k=D,\ldots,1.
  \label{eq:p11-right-contraction-derivative}
\end{aligned}
\]
Thus
\[
  \dot{\widehat z}_t
  =
  \dot R_{t,0}+\dot\tau_t.
  \label{eq:p11-z-derivative}
\]
Equivalently, in integral notation,
\[
  \dot{\widehat z}_t
  =
  2\int \phi_t(r)\dot\phi_t(r)\,\dd\nu(r)+\dot\tau_t.
  \label{eq:p11-z-integral-derivative}
\]

\section{Final Score}

Combining \eqref{eq:p11-scalar} and \eqref{eq:p11-z-derivative}, the
fixed-branch analytical score is
\[
  \boxed{
  \dot{\widehat\ell}_T^{\TT}
  =
  \sum_{t=1}^T
  \left[
  \frac{\dot R_{t,0}+\dot\tau_t}
       {R_{t,0}+\tau_t}
  -
  \dot c_t
  \right].
  }
  \label{eq:p11-final-score}
\]
Every term in \eqref{eq:p11-final-score} is the derivative of the declared
scalar only under the fixed-branch assumptions.  If an implementation uses
the companion code's branch-changing choices, then the derivative is not an
ordinary HMC gradient unless those choices are frozen, smoothed, marginalized,
or moved outside the sampler.

\section{Implementation Checklist}

A future same-scalar prototype must implement the following objects.

\begin{longtable}{@{}p{0.31\linewidth}p{0.58\linewidth}@{}}
\toprule
Object & Required derivative or declaration \\
\midrule
\endhead
Model densities & \(\dot{\log f_\alpha}\), \(\dot{\log g_\alpha}\), and state
gradients if \(A_t,b_t\) are differentiated. \\
Previous TT marginal & \(a_{t-1}\), \(\dot a_{t-1}\), \(\widehat z_{t-1}\),
\(\dot{\widehat z}_{t-1}\). \\
Affine/preconditioner branch & Either fixed numerically, or derivatives
\(\dot A_t,\dot b_t\) and \(\operatorname{tr}(A_t^{-1}\dot A_t)\). \\
Stabilizing constant & Either branch derivative \(\dot c_t\) at the frozen
selected grid point, or a renamed scalar with \(c_t\) fixed numerically. \\
TT core construction & Fixed interpolation or least-squares equations and
\(\dot G_{t,k}\) from \eqref{eq:p11-core-linear-derivative} or
\eqref{eq:p11-core-ls-derivative}. \\
Normalizer & Same value path and derivative path, preferably the direct mass
contraction \eqref{eq:p11-right-contraction}. \\
Branch diagnostics & Fixed ranks, fixed interpolation points, fixed samples,
fixed truncation decisions, fixed QR/SVD signs, and finite positive
normalizers. \\
\bottomrule
\end{longtable}

\section{What Is Not Derived}

This note does not derive the derivative of:
\begin{itemize}[leftmargin=*]
\item changing TT-cross interpolation sets or maxvol/pivot choices;
\item random enrichment samples or debug-sample resampling;
\item rank adaptation, SVD truncation rank changes, or sign flips at repeated
  singular values;
\item ESS-triggered reapproximation and sample doubling;
\item nonunique minimizers in the stabilizing constant \(c_t\);
\item the full MATLAB/Octave companion code as an HMC-ready backend.
\end{itemize}
Those operations can still be valuable for filtering and transport.  They are
not, as stated, a smooth same-scalar gradient path for ordinary HMC.

\section{Minimal Verification Protocol}

For each external parameter coordinate \(\alpha_i\), a future implementation
should run:
\begin{enumerate}[leftmargin=*]
\item value check: the scalar used by HMC is exactly
  \(\sum_t\{\log \widehat z_t-c_t\}\);
\item branch check: ranks, interpolation sets, sample arrays, selected
  constant point, QR/SVD sign convention, and basis/domain maps are unchanged
  under \(\alpha_i\pm h\);
\item finite-difference parity:
  \[
    \frac{\widehat\ell_T^{\TT}(\alpha+h e_i)
          -\widehat\ell_T^{\TT}(\alpha-h e_i)}{2h}
    \approx
    \dot{\widehat\ell}_T^{\TT}(\alpha);
  \]
\item failure check: if branch objects change, report branch-local derivative
  failure rather than tuning HMC on the reported gradient.
\end{enumerate}

\section{Conclusion}

The fixed-branch Zhao--Cui squared-TT scalar has a clear analytical derivative:
differentiate the input density, differentiate the fixed TT core equations,
contract the differentiated cores through the same mass matrices used by the
normalizer, and subtract the derivative of the stabilizing constant used in
the declared scalar.  This is enough for a future implementation plan.  It is
not enough to claim that the adaptive companion code already supplies an HMC
gradient.

\begin{thebibliography}{9}
\bibitem{zhao-cui-2024}
Zhao, B. and Cui, T. (2024).
\newblock Tensor-Train Methods for Sequential State and Parameter Learning in
State-Space Models.
\newblock \emph{Journal of Machine Learning Research}, 25(272):1--83.

\bibitem{cui-dolgov-2022}
Cui, T. and Dolgov, S. (2022).
\newblock Deep Composition of Tensor-Trains Using Squared Inverse Rosenblatt
Transports.
\newblock \emph{Foundations of Computational Mathematics}, 22:1863--1922.
\end{thebibliography}

\end{document}
